% Transcription — fonds Alexandre Grothendieck, shelfmark 63, batch 4 (pages 61-80).
% Established from the facsimile archives/batches/63/batch-04.pdf (a copy of
% 63.pdf fetched through the project relay, no cover sheet: PDF page k =
% archivists' page 60+k), per the skill transcribe-grothendieck. The folder is
% a hundred pages in five batches; this is the fourth. Aligned with the
% decisions of batches 1-3 (transcripts/63/batch-01.fr.tex ... batch-03.fr.tex).
% Folder 62 (transcripts/62/) is the notational reference for the group
% « Courbes elliptiques » (62-64).
%
% Pass: Opus 5.5 (claude-opus-5-5), 2026-09-26 — first pass, unchecked against the pages by a human.
%
% Pages skipped:
% - 64 and 78, blank cream sheets.
% - 67, 69, 71: typed versos of the yellowed sheets carrying pp. 66, 68, 70
%   (Bourbaki drafts, « n° 357 » p. 19 on integration of representations,
%   « n° 380 » pp. 6-7 on normed algebras); nothing in his hand, unrelated.
% - 73 and 77: typed leaves of a draft of the SGA 4 Avant-propos (the
%   bibliography page ending « GT », and typed p. 1 « AVANT-PROPOS »), on
%   the backs of pp. 72 and 76, crossed through by one diagonal line;
%   nothing of his beyond the cancelling stroke, so skipped.
%
% Pages given one \note each, not re-set (someone else's text):
% - Pages 61-63: the end of the typescript copy of A. O. L. Atkin,
%   « Congruence Hecke operators », begun on p. 56 (batch 3): typed p. 6
%   (header cut off; §3 « The subgroup Gamma_0(q) for q prime », Theorem 2),
%   « atkin page 7 » (Conjecture 3) and « atkin page 8 » (the closing
%   discussion and the References 1-9). As in batch 3, every handwritten
%   mark is part of the photocopied image (same grey as the typing): the
%   author's hand-filled Gamma's, sums and Greek letters, a copy-editor's
%   underlinings (wavy = bold, straight = italic) and « # » space marks,
%   and a few English corrections (« for all alpha >= 1 », « k and », « k + »,
%   « are constants » on p. 62; « 136-148 » replacing a page range in
%   reference 5 on p. 63). None is attributed to Grothendieck.
% - Page 75: typed p. « - 3 - » of the same SGA 4 Avant-propos draft (§2 b-c
%   on Igusa, Ogg-Chafarevitch; §3 on exposés I-XIX and footnotes (2),
%   (3)), on the back of p. 74, crossed through by one diagonal line. It
%   carries his ink corrections, so it is kept with \page{75} and one
%   French \note giving them in full with their place (rule used for
%   folder 36's SGA 7 leaves and folder 62's pencilled Bourbaki leaves);
%   the typed text is not re-set.
% - Pages 79-80: a two-leaf typescript (typed page numbers « X-1 », « X-2 »),
%   « Liste des courbes elliptiques à multiplication complexe dont
%   l'invariant modulaire j est rationnel »: the thirteen orders of class
%   number 1 (d, f) and their j, from Weber's Lehrbuch der Algebra III, with
%   « plus peut-être une autre valeur d_? ». Decision: treated as another
%   author's text, not transcribed in full. Evidence: no author is named on
%   the leaves and it is not in his hand; the numbering « X-1 », « X-2 »
%   marks an exposé X of some series; the typewriter differs from Deligne's
%   typescript of pp. 15-33; and the one worded pencil note (« NB les d
%   sont premiers ! », p. 79) is a reader's comment on the text, not an
%   author's revision. Identification of the author is left open. The
%   pencil marks, which look like his fast French hand (attributed with the
%   same caution as batch 2), are given in full as \marginal.
%
% Pages transcribed in full (his own hand, ink): 65, 66, 68, 70, 72, 74, 76
% — seven leaves, none a photocopy, on the periods of an elliptic curve and
% the comparison of the algebraic splitting (omega, eta) of H^1_dR with the
% Hodge splitting (omega, omega-bar).
%
% His pagination: none on these leaves. Numbered runs: Atkin's (H1q)-(H4q),
% Theorem 2, Conjecture 3 (i)-(iv), References 1-9; the typescript's
% entries 1.1, 1.2, 2.1, 3.1-3.3, 7.1, 7.2, 11.1, 19.1, 43.1, 67.1, 163.1,
% d_?.1; his own « Questions a), b) » (p. 66). Dates: none written by him;
% Atkin's references run to 1967; the CM list, which still allows « une
% autre valeur d_? », predates or ignores the Heegner-Stark-Baker
% resolution of the class-number-one problem (1966-67).
%
% What the batch is about: the end of Atkin's typescript; then his own
% notes on the analytic side of the elliptic curve y^2 = x^3 + g4 x + g6:
% the normalisation by a scalar omega (x = omega^-2 x2, ...), omega = dx/2y,
% eta = x omega, periods omega_i and quasi-periods eta_i via the
% Weierstrass series, the passage from the splitting (omega, eta) of the
% forms of the second kind to the Hodge splitting (omega, omega-bar),
% eta + alpha omega = beta omega-bar, the explicit Eisenstein-type series
% for alpha in terms of tau, a linear-algebra account (V = L + L^-1, the
% involution phi, beta + beta-bar = 0), beta = -2i rho with rho = Vol(X),
% 4 rho sigma = 1 - alpha alpha-bar, and the normalisation omega^12 = Delta
% making (arg alpha)^6 a canonical invariant, with the question of
% expressing it through j. Then the CM typescript.
%
% \ill{}: 15. \uncertain{}: 23. Hardest: the prose of pp. 72 and 76 (fast
% hand) and the sums on p. 70 (overwritten numerators).

\documentclass[11pt,a4paper]{article}
% Le préambule commun à toutes les transcriptions du fonds Grothendieck.
%
% Il tient en une idée : **l'incertitude se compose, elle ne se résout pas.**
% Une transcription qui lisse un mot douteux a détruit la seule chose pour
% laquelle on la faisait — un lecteur doit pouvoir savoir, à chaque endroit,
% ce qui était sur le feuillet et ce qui a été ajouté. Les six macros qui
% suivent sont donc l'appareil critique, et non de la décoration.
%
% Les mêmes commandes sont reconnues par `scripts/render.mjs`, qui produit la
% vue de lecture du site. Une macro ajoutée ici doit l'être aussi là-bas,
% faute de quoi le rendu échouera bruyamment — ce qui est voulu : un rendu
% silencieusement faux est pire qu'un rendu absent.

\ProvidesPackage{grothendieck}[2026/08/08 Transcriptions du fonds Grothendieck]

\usepackage{amsmath,amssymb,amsthm}
\usepackage{mathrsfs}
\usepackage{stmaryrd}
% Les diagrammes commutatifs sont partout dans ce fonds : c'est la langue
% même de la géométrie algébrique de Grothendieck, et une transcription qui
% les rendrait en prose ne transcrirait rien.
\usepackage{tikz-cd}
% Français par défaut — c'est la langue des feuillets — mais l'anglais est
% chargé aussi : la traduction et le résumé sont des documents anglais, et la
% typographie française (espaces avant les deux-points) y serait une faute.
% Ces documents basculent par \selectlanguage{english} en tête de corps.
\usepackage[english,french]{babel}
\usepackage{fontspec}
\usepackage{geometry}
\geometry{a4paper,margin=25mm,marginparwidth=28mm}
\usepackage{marginnote}
\usepackage{xcolor}
% ulem, not soul. `soul`'s \st and \ul are built on a character-by-character
% scan that XeLaTeX's font machinery breaks: the document fails with an
% undefined control sequence rather than degrading. ulem does the same two jobs
% and is Unicode-engine safe. `normalem` stops it hijacking \emph, which the
% transcriptions use for Grothendieck's own emphasis.
\usepackage[normalem]{ulem}
\usepackage{hyperref}
\hypersetup{colorlinks=true,linkcolor=black,urlcolor=black,pdfborder={0 0 0}}

% --- Métadonnées du lot -----------------------------------------------------
% Reprises telles quelles par le script de rendu, qui les affiche en tête de
% la vue de lecture. Elles fixent ce que le fichier prétend couvrir : sans
% elles, une transcription flotte sans point d'attache dans un fonds de seize
% mille feuillets.

\newcommand{\folder}[1]{\gdef\grothFolder{#1}}
\newcommand{\batch}[1]{\gdef\grothBatch{#1}}
\newcommand{\leaves}[2]{\gdef\grothFirst{#1}\gdef\grothLast{#2}}
\newcommand{\foldertitle}[1]{\gdef\grothFoldertitle{#1}}
\gdef\grothFolder{?}\gdef\grothBatch{?}\gdef\grothFirst{?}\gdef\grothLast{?}\gdef\grothFoldertitle{}

% --- Le résumé --------------------------------------------------------------
%
% Il ouvre la lecture modernisée, sous le titre « Résumé », et il est composé
% en petit. Ce n'est pas un ornement : c'est le seul passage du document qui
% ne soit pas des mathématiques, et un lecteur doit pouvoir le reconnaître
% avant de l'avoir lu — puis le sauter s'il connaît déjà le sujet, ou s'y
% arrêter s'il ne le connaît pas. Le filet qui le suit marque la reprise du
% corps.

\newenvironment{resume}
  {\small\setlength{\parskip}{0.45em}}
  {\par\medskip\hrule\medskip}

% --- Mots-clés ---------------------------------------------------------------
%
% En anglais, en clôture du résumé de la lecture modernisée : le vocabulaire
% moderne sous lequel chercher ce que les feuillets construisent. Le manifeste
% du site (`npm run manifest`) les extrait pour étiqueter la cote entière —
% cette ligne est la source unique des tags, il n'y a pas de fichier à côté.
\newcommand{\keywords}[1]{\par\smallskip\noindent
  {\footnotesize\textsc{Keywords} --- \textit{#1}}\par}

% --- L'appareil critique ----------------------------------------------------

% Le numéro de feuillet, en marge. C'est l'ancre qui permet de régler un
% désaccord sur une formule en regardant *un* feuillet plutôt qu'une chemise
% de six cents. Le site s'en sert aussi pour tourner les pages du fac-similé
% quand on fait défiler la transcription.
\newcommand{\leaf}[1]{%
  \marginnote{\footnotesize\color{gray!70}\textsf{#1}}%
  \label{leaf:#1}\ignorespaces}

% Illisible : on ne devine pas. Un mot inventé qui se lit comme les autres est
% le pire résultat possible.
\newcommand{\ill}{\textcolor{red!60!black}{[\,\ldots\,]}}

% Lecture douteuse : le mot est donné, et signalé comme incertain.
\newcommand{\uncertain}[1]{\uline{#1}}

% Ajout de l'éditeur — une lettre manquante, un mot sous-entendu.
\newcommand{\add}[1]{\textcolor{blue!50!black}{[#1]}}

% Biffé par Grothendieck lui-même. Ce qu'il a rayé fait partie du document :
% une rature dit souvent où la pensée a changé de direction.
\newcommand{\struck}[1]{\sout{#1}}

% Note du transcripteur, sur l'état matériel du feuillet.
\newcommand{\note}[1]{\par\smallskip{\small\color{gray!60!black}\textit{[#1]}}\par\smallskip}

% Note marginale de Grothendieck, distincte du corps du texte.
\newcommand{\marginal}[1]{\par\smallskip\noindent\hspace{1em}%
  {\small\color{gray!60!black}#1}\par\smallskip}

% --- Titre ------------------------------------------------------------------

\AtBeginDocument{%
  \begin{center}
    {\large\bfseries Fonds Alexandre Grothendieck --- cote n\textsuperscript{o}~\grothFolder}\\[2pt]
    {\itshape\grothFoldertitle}\\[4pt]
    {\small feuillets \grothFirst--\grothLast\ (lot \grothBatch)}\\[2pt]
    {\footnotesize\color{gray!60!black}Transcription établie d'après le fac-similé numérisé
     par l'Université de Montpellier.\\
     Les lectures incertaines sont soulignées, les illisibles notées [\,\ldots\,],
     les ajouts entre crochets.}
  \end{center}
  \vspace{1em}}

\endinput


\folder{63}
\batch{4}
\pages{61}{80}
\dating{[à partir de 1964-vers 1970]}
\watermark{Édition de démonstration}
\foldertitle{Formulaire courbes elliptiques : notes et copies de notes manuscrites (s.d.), tapuscrit (s.d.), copies d'articles (s.d.)}

\begin{document}

\page{61}
\note{Les pages 61 à 63 terminent la copie du tapuscrit de A.~O.~L.~Atkin,
« Congruence Hecke operators », commencée page 56. Texte d'un autre auteur :
il n'est pas recomposé ; chaque page reçoit une note. Comme au lot
précédent, toutes les marques manuscrites font partie de l'image photocopiée
(formules complétées par l'auteur, soulignements et signes d'un correcteur
d'épreuves, corrections en anglais) ; aucune n'est de sa main.}
\note{Page dactylographiée 6 (en-tête coupé) : §3 « The subgroup
$\Gamma_0(q)$ for $q$ prime » : la théorie du §2 s'étend en remplaçant $T_q$
par $U_qF(z)=q^{-1}\sum_{\lambda=0}^{q-1}F\bigl(\frac{z+\lambda}{q}\bigr)
=\sum a(nq)x^n$ ; formes « nouvelles » $C^0$ ; propriétés (H1q)-(H4q) ;
Théorème 2 : base de formes propres simultanées des $T_p$ et de $U_q$,
équivalente aux relations $a(np)-a(n)a(p)+p^{-k-1}a(n/p)=0$ ($p\neq q$) et
$a(nq)-a(n)a(q)=0$ ; pour $k$ négatif, $a(q)=\pm q^{-k/2-1}$ ; renvoi à un
article de J.~Lehner et de l'auteur pour $\Gamma_0(m)$. Marques : le mot
« entire » et « an extension » ajoutés à la machine au-dessus de la ligne,
soulignements ondulés et droits, le signe $\neq$ tracé à la main, un signe
dièse dans la marge droite.}

\page{62}
\note{Page « atkin page 7 » : Conjecture 3 ($q\leq 23$ premier) : il existe
des fonctions entières $F_\alpha(x)=x+\sum_{n\geq 1}a_\alpha(n)x^n$ sur
$\Gamma_0(q)$, nulles en $z=i\infty$, vérifiant (i)
$a_\alpha(np)-a_\alpha(n)a_\alpha(p)+p^{-1}a_\alpha(n/p)\equiv 0$, (ii)
$a_\alpha(nq)-a_\alpha(n)a_\alpha(q)\equiv 0$, (iii)
$a_{\alpha+1}(n)\equiv a_\alpha(n)$ modulo $q^\alpha$, et (iv) toute
fonction entière $F$ à série de Fourier entière vérifie
$U_q^{\alpha+\beta}F(x)\equiv k+k_{\alpha\beta}F_\alpha(x)\pmod{q^\alpha}$ ;
remarques sur l'unicité de $F_\alpha$, l'attrait d'une fonction à
coefficients $q$-adiques, et l'espoir d'une preuve par la géométrie
algébrique. Corrections manuscrites en anglais, dans l'image photocopiée :
« for all $\alpha\geq 1$ » ajouté après « we have », « $k$ and » inséré
devant « $k_{\alpha\beta}$ », « $k+$ » ajouté dans la congruence, « is a
constant » corrigé en « are constants » ; soulignements, un signe dièse.}

\page{63}
\note{Page « atkin page 8 » : extension de la Conjecture 2 comme la
Conjecture 3 étend la Conjecture 1 ; les deux difficultés (l'analogue de
(H4q), qu'on ne sait obtenir que si $\Gamma_0(q)$ est de genre 0, et un
critère de diagonalisation simultanée) ; puis « References » 1 à 9
(Atkin 1967 ; Atkin, Blaricum 1966 ; Atkin et O'Brien 1967 ; Gunning 1962 ;
Lehner 1949, deux articles ; Lehner 1965 ; Rankin 1964 ; Rankin et
Rushforth 1954), et une ligne d'adresse dactylographiée à peine lisible.
Marques : « for all $\alpha$ » ajouté à la machine ; soulignements ondulés
sous les numéros de volume ; dans la référence 5, une pagination biffée et
remplacée à la main par « 136-148 » ; « Modular », « Functions »,
« Forms » corrigés à la main dans la référence 8 ; un signe dièse final.}

\page{65}
\[
x_2,\ y_3,\ g_4,\ g_6 \qquad \boxed{y_3^2 = x_2^3 + g_4x_2 + g_6}
\]
\[
\begin{cases}
x = x_\omega = \omega^{-2}x_2\\
y = y_\omega = \omega^{-3}y_3\\
\gamma_4 = \gamma_{4\omega} = \omega^{-4}g_4\\
\gamma_6 = \gamma_{6\omega} = \omega^{-6}g_6
\end{cases}
\qquad \text{i.e.}\quad y^2 = x^3 + \gamma_4x + \gamma_6
\]
\note{$\gamma_4$ est écrit sur un $g_4$ biffé.}
\[
\boxed{dx = 2y\,\omega}\quad \text{i.e.}\quad
\boxed{d(\omega^{-2}x_2) = 2\omega^{-2}y_3}
\]
\[
2y\,dy = (3x^2+\gamma_4)\,dx = 2y\,(3x^2+\gamma_4)\,\omega
\]
\[
\boxed{dy = (3x^2+\gamma_4)\,\omega}\qquad
\boxed{d(\omega^{-3}y_3) = (3x_2^2+g_4)\,\omega^{-3}}
\]
\note{dans les deux cadres, un coefficient numérique en fraction est
biffé devant la parenthèse, et dans le second deux facteurs devant
$x_2^2$ et $g_4$ sont noircis : les formules se lisent comme données
ici.}
\[
\begin{cases}
\omega = \dfrac12\,\dfrac{dx}{y}\\[2mm]
\eta = \omega x = \dfrac12\,\dfrac{x\,dx}{y}
\end{cases}
\qquad [\,= \omega^{-1}x_2\,]
\]
\note{un mot est noirci après $\omega$ à la première ligne, et un autre
avant $\frac12$ à la seconde.}
\[
\begin{cases}
\int_{\gamma_1}\omega = \omega_1, & \omega_1 = \int_0^{\omega_1}dz\\[1mm]
\int_{\gamma_2}\omega = \omega_2, & \omega_2 = \int_0^{\omega_2}dz
\end{cases}
\qquad X^{\mathrm{an}} \simeq \mathbf{C}/E(\omega_1,\omega_2)
\]
\[
\begin{cases}
x = \dfrac{1}{z^2} + \sum'_{a\in E(\omega_1,\omega_2)}
\Bigl[\dfrac{1}{(z-a)^2} - \dfrac{1}{a^2}\Bigr]\\[3mm]
y = -\dfrac{1}{z^3} - \sum_{a\in E(\omega_1,\omega_2)}\dfrac{1}{(z-a)^3}
\end{cases}
\]
\note{dans la formule de $y$, le signe devant la seconde somme est
surchargé ; nous lisons $-$, sans certitude.}
\[
\begin{cases}
\omega = dz\\
\eta = x\,dz
\end{cases}
\qquad \eta_1 = \int_u^{u+\omega_1}x\,dz,\quad
\eta_2 = \int_u^{u+\omega_2}x\,dz
\]
\[
\gamma_4 = \qquad \gamma_6 =
\]
\note{les deux dernières lignes sont laissées sans second membre.}

\page{66}
Passage du splitting standard des formes différentielles de 2$^{\mathrm{e}}$
espèce en $\omega,\eta$ au splitting transcendant (de Hodge) $\omega,\bar\omega$,
donc on aura
\[
\begin{cases}
\eta + \alpha\omega = \beta\bar\omega\\[1mm]
\alpha = -\dfrac{\bar\omega_1\eta_2 - \bar\omega_2\eta_1}
{\bar\omega_1\omega_2 - \bar\omega_2\omega_1}
\end{cases}
\]
\marginal{NB. dénominateur $= 2i\,|\omega_1|^2\,\mathrm{Im}\,\tau \neq 0$}

\emph{NB} On a une 2-forme invariante canonique associée à la comparaison
des deux splittings
\[
\pi = \alpha\,\omega^2 \in \Gamma\,\mathrm{Hom}(\underline{\omega}^{-1},
\underline{\omega}) = \Gamma\,\underline{\omega}^2
\]

\emph{Questions} a) Quand $\alpha=0$ ?, i.e.\ $\bar\omega$ proportionnel
à $\eta$ ??

b) Quand $\alpha$ est-il algébrique sur $\mathbf{Q}$ ?

\page{68}
\[
\eta_1 = \Bigl[-\frac1z\Bigr]_u^{u+\omega_1}
+ \sum_{a\neq 0}\Bigl[-\frac{1}{z-a} - \frac{1}{a^2}\,z\Bigr]_u^{u+\omega_1}
\]
\[
= \Bigl(\frac1u - \frac{1}{u+\omega_1}\Bigr)
+ \sum_{a\neq 0}\Bigl[\frac{1}{u-a} - \frac{1}{u+\omega_1-a}
- \frac{\omega_1}{a^2}\Bigr]
\]
\[
= \omega_1\Bigl[\frac{1}{u(u+\omega_1)}
+ \sum_{a\neq 0}\Bigl(\frac{1}{(u-a)(u+\omega_1-a)} - \frac{1}{a^2}\Bigr)\Bigr]
\]
\note{sous le terme de la somme, une accolade donne sa valeur :
$\frac{1}{a^3}\,\frac{(2u+\omega_1)-u(u+\omega_1)a^{-1}}
{(1-ua^{-1})(1-(u+\omega_1)a^{-1})}$ ; une ligne commencée
« $=\omega$ » est biffée.}
\[
= \omega_1\Bigl[\frac{1}{u(u+\omega_1)} + \frac{1}{(u-\omega_1)u}
- \frac{1}{\omega_1^2} + \sum_{\substack{a\neq 0\\ a\neq\omega_1}}\cdots\Bigr]
\]
\note{une flèche renvoie la somme au terme de l'accolade précédente ; sous
les deux premiers termes, une accolade :
$\frac{2}{(u+\omega_1)(u-\omega_1)}$, un facteur $u$ biffé au
dénominateur.}

faisons $u\to 0$, on trouve
\[
\eta_1 = \omega_1\Bigl[-\frac{3}{\omega_1^2}
+ \sum_{\substack{a\neq 0\\ a\neq\omega_1}}\frac{1}{a^3}\,
\frac{\omega_1}{1-\omega_1a^{-1}}\Bigr]
\]
\note{accolade sous le terme : $\frac{1}{a^3}\,\frac{1}{\omega_1^{-1}-a^{-1}}$.}
\[
\eta_1 = -\frac{3}{\omega_1} + \omega_1\sum_{\substack{a\neq 0\\ a\neq\omega_1}}
\Bigl(\frac{1}{a^3}\,\frac{1}{\omega_1^{-1}-a^{-1}}\Bigr)
= -\frac{3}{\omega_1} + \omega_1^2\sum_{\substack{a\neq 0\\ a\neq\omega_1}}
\frac{1}{a^2}\cdot\frac{1}{a-\omega_1}
\]
\note{accolade sous le terme du milieu : $\frac{1}{a^2}\,\frac{\omega_1}{a-\omega_1}$.}
\[
a = m\omega_1 + n\omega_2
\]
\[
\eta_2 = \qquad -\frac{3}{\omega_2} + \omega_2^2\sum_{\substack{a\neq 0\\ a\neq\omega_2}}
\frac{1}{a^2}\,\frac{1}{a-\omega_2}
\]
\note{la formule de $\eta_2$ est écrite à droite, sous celle de $\eta_1$ ;
l'indice de la somme se lit « $a\neq\omega_1$ », par report de la ligne
précédente, sans doute pour $a\neq\omega_2$. En bas à gauche, un mot biffé
\ill{}.}

\page{70}
\[
\eta_1 = -\frac{3}{\omega_1} + \omega_1^2
\sum_{\substack{(m,n)\neq(0,0)\\ (m,n)\neq(1,0)}}
\frac{1}{(m\omega_1+n\omega_2)^2}\,\frac{1}{(m-1)\omega_1+n\omega_2}
\]
\[
\eta_2 = -\frac{3}{\omega_2} + \omega_2^2
\sum_{\substack{(m,n)\neq(0,0)\\ (m,n)\neq(0,1)}}
\frac{1}{(m\omega_1+n\omega_2)^2}\,\frac{1}{m\omega_1+(n-1)\omega_2}
\]
\[
\bar\omega_1\eta_2 - \bar\omega_2\eta_1
= -3\Bigl(\frac{\bar\omega_1}{\omega_2} - \frac{\bar\omega_2}{\omega_1}\Bigr)
+ \sum_{(m,n)\neq(0,0),(1,0),(0,1)} \frac{N_{m,n}}{D_{m,n}}
\]
\[
N_{m,n} = |\omega_1|^2\Bigl((m-1)\omega_2^2 + n\,\frac{\omega_2^3}{\omega_1}\Bigr)
- |\omega_2|^2\Bigl((n-1)\omega_1^2 + m\,\frac{\omega_1^3}{\omega_2}\Bigr)
\]
\[
D_{m,n} = (m\omega_1+n\omega_2)^2\,\bigl((m-1)\omega_1+n\omega_2\bigr)
\bigl(m\omega_1+(n-1)\omega_2\bigr)
\]
\[
+\ \bar\omega_1\omega_2^2\,\frac{1}{\omega_1^2(\omega_1-\omega_2)}
- \bar\omega_2\omega_1^2\,\frac{1}{\omega_2^2(\omega_2-\omega_1)}
\]
\note{la notation $N_{m,n}/D_{m,n}$ est nôtre : il écrit la fraction d'un
seul tenant, sur deux lignes. Un premier essai de somme après la
parenthèse, avec les exclusions $(0,0)$, $(1,0)$, $(0,1)$, est biffé de
hachures. Sous les deux derniers termes (les termes $(0,1)$ de $\eta_1$ et
$(1,0)$ de $\eta_2$), une accolade :
$\frac{\bar\omega_1\omega_2^4+\bar\omega_2\omega_1^4}{(\omega_1\omega_2)^2(\omega_1-\omega_2)}$,
un exposant biffé sur $(\omega_1-\omega_2)$.}
\[
= \frac{\bar\omega_1\omega_2^4+\bar\omega_2\omega_1^4}{\omega_1^2\omega_2^2(\omega_1-\omega_2)}
- 3\,\frac{|\omega_1|^2-|\omega_2|^2}{\omega_1\omega_2}
+ \sum_{(m,n)\neq(0,0),(1,0),(0,1)}\cdots
\]
\note{une flèche renvoie la somme à celle de la ligne précédente.}

\[
= \frac{|\omega_1|^2}{\omega_1^2}\,\frac{\tau^4+\bar\tau}{\tau^2(1-\tau)}
- \frac{|\omega_1|^2}{\omega_1^2}\,\frac{1-|\tau|^2}{\tau}
\]
\[
+ \sum_{(m,n)\neq(0,0),(0,1),(1,0)}\frac{|\omega_1|^2}{\omega_1^2}\,
\frac{(m-1)\tau^2+n\tau^3-(n-1)|\tau|^2+m|\tau|^2\tau^{-1}}
{(m+n\tau)^2\bigl((m-1)+n\tau\bigr)\bigl(m+(n-1)\tau\bigr)}
\]
\note{ici $\tau=\omega_2/\omega_1$, ce que la page ne dit pas. Le premier facteur
est récrit sur une première version noircie, et l'exposant du dénominateur
$\tau^2(1-\tau)$ corrigé ; le dernier terme du numérateur de la somme
($m|\tau|^2\tau^{-1}$) est lu ainsi, sans certitude.}
\[
= \frac{|\omega_1|^2}{\omega_1^2}\Bigl(\frac{\tau^5+|\tau|^2}{\tau^3(1-\tau)}
- \frac{1-|\tau|^2}{\tau}
+ \sum_{(m,n)\neq(0,0),(0,1),(1,0)}
\frac{n\tau^4+(m-1)\tau^3+(n-1)|\tau|^2\tau+m|\tau|^2}
{\tau(m+n\tau)^2\bigl((m-1)+n\tau\bigr)\bigl(m+(n-1)\tau\bigr)}\Bigr)
\]
\note{le signe devant $(n-1)|\tau|^2\tau$ est lu $+$ ; la ligne précédente
porte $-(n-1)|\tau|^2$.}
\[
\alpha = -\frac{\bar\omega_1\eta_2-\bar\omega_2\eta_1}{2i\,|\omega_1|^2\,\mathrm{Im}\,\tau}
= -\frac{1}{2i\,\omega_1^2\,\mathrm{Im}(\tau)}
\Bigl(\frac{\tau^5+|\tau|^2}{\tau^3(1-\tau)} - \frac{1-|\tau|^2}{\tau}
+ \sum_{(m,n)\neq(0,0),(0,1),(1,0)}\cdots\Bigr)
\]
\note{un $\bar\omega_1$ biffé précède $\alpha$ ; une flèche renvoie la
somme à la ligne précédente.}

\page{72}
\[
V = L + L^{-1}
\]
$L$ défini par $\mathbf{T}$ sur $\mathcal{U}=\lbrace z,\ |z|=1\rbrace$,
$P$ sur $\mathbf{R}^{*+}$

$\varphi\colon \bar V \xrightarrow{\ \sim\ } V$ involution
\[
\bar L \to L + L^{-1} \qquad \alpha\in \bar L^{-1}\otimes L,\quad
\beta\in(\bar L L)^{-1}
\]
\note{sous $\bar L^{-1}\otimes L$ : « défini par $T^{(2)}$ sur
$\mathcal{U}$ » ; sous $\bar L L$ : « défini par $P^{(2)}$ sur
$\mathbf{R}^{*+}$ », et au-dessus de la ligne suivante, entre parenthèses,
« (\ill{} \uncertain{ailleurs arbitraires}) ».}

$\boxed{\beta\neq 0}$ \quad \uncertain{Alors} \uncertain{les conditions}
\ill{} $\alpha$ et $\beta$ \uncertain{découlent} \ill{} l'involution \ill{}
\ill{}.
\[
V = L + \varphi(L)
\]
H \uncertain{dual} de $V$ \uncertain{s'écrit} $z = x + \bar y$
\struck{\ill{}} \ ($x\in L$, $y\in L$ \uncertain{uniques})
\[
\alpha\,\bar z = y + \bar x
\]
\[
\varphi(x) = \bar x = \alpha(\tilde x) + \beta(\tilde x)
\]
\note{un mot biffé à droite de cette formule ; les tildes sur $x$ sont
lus ainsi, sans certitude.}

si $y\in L^{-1}$, on aura
\[
y = a + \varphi(b) \quad a\in L,\ b\in L;\qquad
= \bigl(a+\alpha(b)\bigr) + \beta(b)
\]
donc $a+\alpha(b)=0$, $\beta(b)=y$, d'où $b=\beta^{-1}y$,
$a=-\alpha(b)=-\alpha\beta^{-1}(y)$

donc
\[
y = -\alpha\beta^{-1}(y) + \overline{\beta^{-1}(y)}, \quad\text{d'où}
\]
\[
\bar y = \beta^{-1}(y) - \overline{\alpha\beta^{-1}(y)}
= \bigl(\beta^{-1}(y) - \alpha^2\beta^{-1}(y)\bigr) - \beta\alpha\beta^{-1}(y)
\]
\note{les deux lignes sont très surchargées : un « $+\alpha$ » biffé dans
la première, un « $\varphi(y)$ » biffé au-dessus de $\bar y$, des signes
noircis et un facteur $\alpha$ biffé devant $\beta^{-1}(y)$ dans la
seconde. Nous donnons l'état final lisible.}

Sous quelle condition, pour la forme \add{\uncertain{alternée}}
\uncertain{canonique} sur $V$, a-t-on
$\langle\bar x,\bar y\rangle = \overline{\langle x,y\rangle}$ ? Il suffit de
l'écrire pour $x\in L$, $y\in L^{-1}$, on a alors
\[
\langle\bar x,\bar y\rangle = \langle\underbrace{\alpha(x)}_{L}
+ \underbrace{\beta(x)}_{L^{-1}},\ \bigl(\beta^{-1}(y)-\alpha^2\beta^{-1}(y)\bigr)
- \beta\alpha\beta^{-1}(y)\rangle
\]
\[
= -\langle\alpha(x),\beta\alpha\beta^{-1}(y)\rangle
+ \langle\beta^{-1}y-\alpha^2\beta^{-1}y,\ \beta(x)\rangle
\]
\emph{On trouve la condition} : $\beta+\bar\beta=0$ \emph{i.e.\ $\beta$
purement imaginaire}.

\page{74}
\struck{\ill{}}

Or dans le cas qui nous occupe, on aura (si $\omega$ est une 1-forme
invariante)
\[
\beta(\omega,\omega) = \int_X\omega\bar\omega = \int_X(dx+i\,dy)(dx-i\,dy)
= -2i\int_X dx\wedge dy = -2i\,|\omega_1|^2\,\mathrm{Im}\,\tau
= |\omega_1|^2(\bar\tau-\tau)
\]
\note{sous l'intégrale du milieu : $x=\mathrm{Re}\,z$, $y=\mathrm{Im}\,z$.}

La condition \uncertain{voulue} exige \uncertain{donc} que
\[
\beta = -2i\rho \quad\text{avec}\quad \rho>0, \qquad \rho = \int_X dx\wedge dy
\]
\note{un premier terme est biffé après « $\beta=$ ».}

Quant à $\alpha$, il vient a priori une classe arbitraire de
$\bar L^{-1}\otimes L$. \uncertain{Notons} que « $\alpha$ » est
(indépendamment du choix d'une base $\omega$ de $L=\underline{\omega}$)
défini, qu'est-il ?

[NB. $\gamma = -\alpha\beta^{-1}\in L^{\otimes 2}$ est l'objet qui fait
passer du splitting « algébrique » de $V$ au « splitting transcendant »
\add{de $V$} (i.e.\ de $x\mapsto\bar x$). La connaissance de $(\alpha,\beta)$
équivaut à celle de $\beta$ et $\gamma$, i.e.\ de la \uncertain{forme
volume} $\int_X\omega\bar\omega$ et des splittings.]

Posons, si $\omega,\eta$ sont choisis [donc $L\simeq\mathbf{C}$,
$L^{-1}\simeq\mathbf{C}$ isomorphismes donnés], en plus de
\struck{$\rho = \mathrm{Vol}(X)$},
\[
\rho = -\frac{1}{2i}\int\omega\bar\omega = \mathrm{Vol}(X),
\]
\uncertain{la forme}
\[
\sigma = -\frac{1}{2i}\int\eta_0\bar\eta_0
\]
où $\eta_0$ est une forme différentielle \uncertain{fermée} dans la classe
de celle de $\eta$.

On trouve
\[
4\rho\sigma = (1-\alpha\bar\alpha)
\]

\page{75}
\note{Feuillet dactylographié « - 3 - » d'un projet de l'« Avant-propos »
de SGA~4, barré d'un seul trait oblique, au dos de la page 74 : fin du §2
(b) les travaux d'Igusa sur les « vanishing cycles » et l'inégalité de
Picard ; c) les résultats d'Ogg et Chafarevitch) et §3 (plan des exposés I
à VI, VII-VIII, IX-X, XII à XVI, XVII à XIX), avec les notes (2) « Obtenus
entre Septembre 1962 et Mars 1963 par M.~Artin et A.~Grothendieck » et (3)
« Obtenus en Février-Mars 1963 par A.~Grothendieck […] ». Texte non
recomposé. Ses corrections à l'encre, dans l'ordre de la page : au §2 c),
après « à coefficient », « des » surchargé et un mot ajouté au-dessus,
\uncertain{dans} ; des virgules ajoutées après « algébriques » (même
alinéa), après « La présentation » et après « contenus dans Artin », et
après « coefficients de torsion » ; une croix dans la marge gauche en face
de « contenus dans Artin » ; les numéros « VII et VIII » et « IX, X »
repassés ou surchargés à la main ; en face de l'alinéa « Les résultats des
exposés XII à XVI », dans la marge gauche, « XII » souligné, et plus à
gauche, au crayon, « X.3.10 » et « X.3.4 » ; après « Théorèmes de
changement de base" », « Nous avons pris soin de donner » écrit puis biffé,
remplacé, à l'encre bleue sous la ligne, par « qui techniquement
constituent le résultat central de ce séminaire. », un trait ondulé menant
à la marge droite ; dans « Les exposés XVII à XIX », le dernier numéro
surchargé ; « théorème des dualités globales » ramené à « théorème de
dualité globale » par suppression des trois « s » ; dans la note (3),
« dans 1965 » biffé après « qui sera exposé », et « SGA 1965 » écrit à la
suite.}

\page{76}
Ainsi
\[
0\leq\alpha\bar\alpha\leq 1
\]
et de façon précise
\[
1-\alpha\bar\alpha = 4\rho\sigma, \quad\text{i.e.}\quad
\alpha\bar\alpha = 1-4\rho\sigma
\]
Ainsi, $\beta$ et $\alpha\bar\alpha$ sont donnés en termes
\uncertain{volumétriques}. Il reste à donner (pour $\alpha\bar\alpha\neq 0$)
$\arg\alpha$ (qui dépend bien \ill{} \ill{} du choix de $\omega$).
Normalisons $\omega$ de façon que $\omega^{12}=\Delta$, ce qui le détermine
\uncertain{mod} multiplication par $\zeta$, $\zeta^{12}=1$, on \uncertain{voit}
que $\arg\alpha\in\mathcal{U}=\lbrace z\in\mathbf{C},\ |z|=1\rbrace$ se trouve
déterminé à \struck{puissance} mult.\ par $\zeta$, avec $\zeta^6=1$.
\uncertain{Donc} \ill{} \ill{} un invariant canonique, (\ill{}
\uncertain{comment} \ill{} en fonction de $j$ !)
\[
(\arg\alpha)^6\in\mathcal{U}.
\]
NB. $\rho$ et $\sigma$ sont déterminés \uncertain{mod rien du tout} par la
normalisation $\omega^{12}=\Delta$).

\page{79}
\note{Les pages 79 et 80 sont un tapuscrit de deux feuillets, paginés
« X-1 » et « X-2 », « Liste des courbes elliptiques à multiplication
complexe dont l'invariant modulaire $j$ est rationnel ». L'auteur n'est pas
nommé et le texte n'est pas de sa main ; la pagination est celle d'un
exposé X d'une série ; la seule note rédigée au crayon est un commentaire
de lecteur. Nous le traitons comme le texte d'un autre auteur : il n'est
pas recomposé, chaque page reçoit une note, et les marques au crayon, d'une
écriture qui paraît être la sienne (attribution incertaine, comme au lot 2),
sont données en entier.}
\note{Feuillet X-1 : ces courbes s'obtiennent à partir des ordres de corps
quadratiques imaginaires de nombre de classes 1 : a) les corps
$\mathbf{Q}(\sqrt{-d})$ d'anneau des entiers principal,
$d=1,2,3,7,11,19,43,67,163$ « plus peut-être une autre valeur $d_?$ » ;
b) les conducteurs $f$ donnant un ordre de nombre de classes 1 : $f=1,2$
pour $d=1$, $f=1,2,3$ pour $d=3$, $f=1,2$ pour $d=7$, $f=1$ pour les autres
(et $d_?$ s'il existe) ; c) pour chaque $(d,f)$, $j(t)$ pour une base
$(1,t)$ de l'ordre, valeurs tirées de H.~Weber, \emph{Lehrbuch der
Algebra}, t.~III, par $j=(f^{24}-16)^3/f^{24}=(f_1^{24}+16)^3/f_1^{24}$ ;
1.1 $(d=1,f=1)$ : $j(i)=2^63^3$ (groupe d'automorphismes d'ordre 4) ; 1.2
$(d=1,f=2)$ : $j(2i)=(2\cdot3\cdot11)^3$. Au crayon : une accolade sur
$d=1,2,3,7,11$ surmontée de « \uncertain{euclidiens} » entre guillemets, et
dans la marge gauche :}
\marginal{NB les $d$ sont premiers !}

\page{80}
\note{Feuillet X-2, suite de la liste : 2.1 $j(\sqrt{-2})=(2^2\cdot5)^3$ ;
3.1 $j\bigl(\frac{-1+\sqrt{-3}}{2}\bigr)=0$ (groupe d'automorphismes
d'ordre 6) ; 3.2 $j(\sqrt{-3})=2^43^35^3$ ; 3.3
$j\bigl(\frac{-1+3\sqrt{-3}}{2}\bigr)=-3\cdot2^{15}5^3$ ; 7.1
$j\bigl(\frac{-1+\sqrt{-7}}{2}\bigr)=-3^35^3$ ; 7.2
$j(\sqrt{-7})=(3\cdot5\cdot17)^3$ ; 11.1
$j\bigl(\frac{-1+\sqrt{-11}}{2}\bigr)=-2^{15}$ ; 19.1 $-(2^5\cdot3)^3$ ;
43.1 $-(2^6\cdot3\cdot5)^3$ ; 67.1 $-(2^5\cdot3\cdot5\cdot11)^3$ ; 163.1
$-(2^6\cdot3\cdot5\cdot23\cdot29)^3$ ; $d_?.1$ : $j=\,?$ ; renvois à Weber
(p.~721, 462, 460, 475). Le texte s'achève par « *** ». Aucune marque
manuscrite, hormis une surcharge dactylographique ($d=3$ en 3.3).}

\end{document}
